% ch01 Synopsis (书页1-6 / p019-024)
% 内容核对说明：原书第1章（书页1--4，p019--p022）为全书Synopsis，纯散文，
% 全章无编号公式、无图表（已对照PNG逐页核验，文字层与MinerU md一致）。
% 书页5（p023）为"Part I"扉页、书页6（p024）为空白页，均不属于本章正文。
% 故本章无 \tag 公式，文末 %TAGS: 清单为空。
\chapter{概要}

原书第1章是一份四页的概要（Synopsis）：不含一个公式、一幅图，任务是画出
全书的论证地图。Baaquie先申明全书的两条主线，澄清"量子"一词的准确含义；
然后按三个组成部分逐一预告第2--10章——每章要解决什么问题、动用什么数学
工具、得到哪一类结果；章末一段说明全书体例：数学细节放进各章附录，可跳过
而不伤连贯性。

对本导读的读者，第1章的正确用法是把它当作"支票清单"：这里点到的每一个
概念——鞅条件、障碍期权、心理未来时间、刚性模型、重正化——都会在后面某章
兑现为精确定义与推导。本章没有任何编号公式，导读的任务相应变为：把每句
预告"翻译"成它在前方章节中的数学落点，并注明"见第N章"。阅读本章不需要
任何先修知识；相反，它是为全书准备的先修地图。

\section{两条主线与"量子"的含义}

本节要解决的问题是全书的名分：一本讲期权与利率的书，凭什么叫"量子"？
作者的回答分两层——两条主线，与一条概念边界。

全书的两条主线。其一：在量子理论的概念与数学框架内定义并分析量化金融
（quantitative finance），特别强调其路径积分（path integral）表述；其二：
把量子场论（quantum field theory）的技术与方法论引入利率（interest rate）
研究。两条主线共用同一套工具，差别在对象的规模：前者（股票期权）是有限
自由度问题，后者（利率曲线）是无穷自由度问题——这条"自由度升级"的路线
正是原书三部分划分的依据，见下面三节的解读。

"量子"的边界。本书不打算用量子理论去改造金融学的基本原理：无套利、货币
的时间价值这些金融学公理原样保留。"量子"一词专指量子理论的抽象数学
构造——概率论、态空间（state space）、算符（operator）、哈密顿量
（Hamiltonian）、对易方程（commutation equations）、拉格朗日量
（Lagrangian）、路径积分、量子化场（quantized field）、玻色子（boson）、
费米子（fermion）等。作者的立场可以概括为：取量子理论的数学，不取其
物理。这对两类读者各是一条信息：物理读者不必寻找波函数的物理诠释；金融
读者可以放心，定价逻辑仍从金融公理出发（见第2--3章）。

为什么路径积分与哈密顿量唱主角。原书给出三层理由。第一，随机量子过程的
路径积分与哈密顿量表述，同随机微积分（stochastic calculus）既等价又相互
独立，而随机微积分正是当今数理金融的基石之一。等价，保证旧结果都能在
新语言中复现——Black--Scholes方程与Merton--Garman方程都将改写为哈密顿量
与路径积分形式（见第4、5章）；独立，保证新语言能走出旧框架——把整条远期
利率曲线当作二维量子场建模，是随机微积分做不到的事（见第7章）。第二，
路径积分为重正化（renormalization）提供了高效的实施框架，而非线性量子
场论恰恰需要重正化才有适定的量（非线性远期利率模型将用到，见第7章）。
第三层隐含在书名的定义里：量子金融（Quantum Finance）指量子理论的概念、
方法与数学同理论及应用金融学的合成——本书即以此合成物为名。
\footnote{对金融背景的读者，"等价"可以落到实处：马尔可夫过程的路径积分
表示与Kolmogorov前向、后向方程描述同一对象，费曼--卡茨（Feynman--Kac）
公式是两种语言之间的桥，第5章将用它把期权定价写成路径积分。}

\section{写作体例：先改写已知，再推导新知}

本节交代全书的教学法承诺，读者可以据此规划自己的阅读路径。

从已知到新知。为了减轻进入量子数学（尤其是路径积分）的门槛，凡引入新
思想，原书都先做"改写"再做"推导"：把金融中已经确立的模型与已知结果用量子
形式重新表述一遍——这是较容易的练习——然后再去推导新结果。作者指出这一
做法的一个意外好处：习惯于随机微积分与偏微分方程（partial differential
equation）的金融理论工作者，看到的是一套与随机微积分完全平行、互为镜像的
形式论，向量子场论数学的过渡因而是平滑的。

自足性与范例。所有重要方程都从金融学第一性原理推导，主要结果尽可能给出
完整而自足的数学处理。若干精确可解模型被反复精讲，充当演示量子金融一般
原理的范例。这本"范例清单"后文可以点名：障碍期权的哈密顿量精确解
（见第4章）、随机波动率期权的路径积分精确解（见第5章）、Vasicek模型的
路径积分解（见第6章）、"刚性"高斯场论的市场拟合（见第7、8章）。这些模型
本身就有金融学上的兴趣，为路径积分的应用提供了现实语境。

数值算法的边界（原书脚注1）。金融问题的路径积分表述为强力计算算法打开了
大门，但数值算法是专门课题，原书除第5.16节一笔带过的蒙特卡罗（Monte
Carlo）算法外不作讨论。

\section{第一部分：金融学基本概念（原书第2--3章）}

原书把全书分为三个组成部分；本节起依次对应。第一部分题为"金融学基本
概念"（Fundamental concepts of finance），回顾金融学的标准概念与期权
理论的方程，任务是为后两个部分定义框架与语境。这部分材料是标准的——
熟悉通行期权教材的读者可以快读。

第2章"金融学导论"（Introduction to finance）面向不熟悉金融的读者，引入
理解本书方程所需的概念与术语：风险与收益（risk and return）、货币的时间
价值（time value of money）、套利（arbitrage）、对冲（hedging），最后是
国债（Treasury Bond）与固定收益证券（fixed income securities）——最后
两项为第二、三部分的利率理论预作铺垫（见第2章）。

第3章"衍生证券"（Derivative securities）引入金融衍生品（financial
derivatives）概念并讨论其定价：回顾经典分析，简述随机微积分的数学，
阐明随机过程与朗之万方程（Langevin equation）的联系，并从第一性原理推导
随机波动率（stochastic volatility）股票期权的价格，即Merton--Garman方程
及其解（见第3章）。值得注意的是，这块"标准材料"的最后一项——波动率本身
的随机性——已经埋下全书的伏笔：让波动率也变成随机对象，正是后文场论思想
登场的入口之一。

\section{第二部分：有限自由度系统（原书第4--6章）}

第二部分题为"有限自由度系统"（Systems with finite number of degrees of
freedom），把哈密顿量与路径积分用于股票期权与随机利率模型。这里出现全书
第一个关键定义：自由度（degree of freedom）指每一时刻$t$的独立随机变量
（random variable）个数。随机演化的股票价格$S(t)$有一个自由度，即期利率
（spot interest rate）$r(t)$同样有一个。这个定义把"系统规模"归结为"每个
时刻随机源的维数"，立刻给出一条分界线：有限自由度系统算出的所有量完全
有限，不需要重正化即适定；只有当自由度成为连续统（第三部分的利率曲线）
时，重正化这类场论手段才成为必需。
\footnote{用随机微分方程的语言说，自由度就是驱动方程的独立布朗运动
（Brownian motion）的条数：$S(t)$或$r(t)$由一条布朗运动驱动，是单自由度
系统；第7章的远期利率曲线则由"每个未来时刻各一条"的布朗运动驱动，是无穷
自由度系统。}

第4章"哈密顿量与股票期权"（Hamiltonians and stock options）把衍生证券
定价整体改写为量子力学问题，四件事各有落点：对常数波动率与随机波动率两种
情形，推导驱动期权价格的哈密顿量（$H_{BS}$与$H_{MG}$，见第4章）；把风险
中性（risk-neutral）演化所需的鞅条件（martingale condition）重新表述为
哈密顿量的语言（见第4章）；证明哈密顿量中的势能项代表一类路径依赖期权
（path-dependent option）（见第4章）；并用相应的哈密顿量精确求解障碍
期权（barrier option）（见第4章）。

第5章"路径积分与股票期权"（Path integrals and stock options）把期权定价
表示为费曼路径积分（Feynman path integral）：第4章的哈密顿量在"期权定价
的偏微分方程"与"它的路径积分实现"之间架起桥梁（见第5章）；书中显式算出
若干路径积分，用来演示路径积分的数学；随机波动率情形被精确解出——这是个
非平凡问题，是演示路径积分精妙之处的范例（见第5章）。此外，随机波动率
路径积分表示带来的精确简化导出一个高效算法，原书借它展示路径积分的数值
方面——即上一节提到的第5.16节蒙特卡罗算法（见第5章）。

第6章"随机利率的哈密顿量与路径积分"（Stochastic interest rates'
Hamiltonians and path integrals）回顾即期与远期利率的重要现有随机模型：
对即期利率得到福克--普朗克（Fokker--Planck）哈密顿量与路径积分，并给出
Vasicek模型的路径积分解（见第6章）；再把Heath--Jarrow--Morton（HJM）
远期利率模型改写为路径积分问题，用路径积分重新推导HJM的已知结果
（见第6章）。原书明言：第6章是为全书的主攻方向——用量子场论建模利率——
所作的准备。

\section{第三部分：利率模型的量子场论（原书第7--10章）}

第三部分题为"利率模型的量子场论"（Quantum field theory of interest rates
models），是本书的核心贡献所在。量子场论是研究无穷多自由度系统的数学
结构；无穷维带来许多随机微积分之外的新特征，最重要的当属非线性场论
（nonlinear field theory）的重正化。本部分所有章节都把远期利率（forward
interest rate）当作量子场处理。

用一行记号即可把"曲线即场"说清楚。远期利率$f(t,x)$是$t$时刻约定的、未来
时刻$x>t$的瞬时贷款利率；固定$t$，曲线
\begin{equation*}
\theta=x-t,\qquad x\mapsto f(t,t+\theta),\qquad \theta\in[0,\infty),
\end{equation*}
就是以未来时间（future time）$\theta$为坐标的一条利率期限结构曲线；第7章
让它随日历时间$t$随机涨落——曲线上每个点各有一个随机源，自由度成为
连续统，"场"由此得名。（记号预告：远期利率速度量子场$\mathcal{A}(t,x)$、
传播子$D(\theta,\theta')$、刚性（rigidity）参数$\mu$与刚度（stiffness）
参数$\lambda$，定义见第7章与卷末记号总表。）

第7章"远期利率的量子场论"（Quantum field theory of forward interest
rates）把路径积分形式用于随机演化曲线：远期利率被建模为随机涨落的曲线，
自然地由量子场论描述。各种线性（高斯）模型（linear (Gaussian) models）
被逐一研究，以展示场论进路的理论灵活性；心理未来时间（psychological
future time）为高斯场论模型提供自然推广（见第7章）；高斯模型的鞅条件被
解出，并给出计价物变换（change of numeraire）的场论推导（见第7章）；
非线性场论则在为取正值的远期利率以及带随机波动率的远期利率建模时自然
出现（见第7章）。

第8章"远期利率的实证与场论模型"（Empirical forward interest rates and
field theory models）转向实证（empiricism）：讨论远期利率的实证性质，
展示如何用欧洲美元（Eurodollar）期货合约（futures contract）的市场数据
校准（calibration）并检验场论模型。本章最重要的结果：所谓"刚性的"
（stiff）高斯场论模型对远期利率的市场行为给出几乎精确的拟合；这项实证
研究为场论刻画远期利率行为的有效性提供了有说服力的证据（见第8章；刚性、
刚度参数的数学身份见第7章）。

第9章"国债衍生品与对冲的场论"（Field theory of Treasury Bonds'
derivatives and hedging）研究国债期货、债券期权（bond option）与利率上限
期权（interest rate cap）的定价；讨论场论利率模型中国债的对冲（hedging），
并证明它是更标准进路的推广；对瞬时与有限时间两种情形都导出精确的对冲
结果，并作半实证（semi-empirical）分析（见第9章）。

第10章"远期利率的场论哈密顿量"（Field theory Hamiltonian of the forward
interest rates）为线性与非线性理论推导态空间（state space）与哈密顿量：
哈密顿量表述给出非线性远期利率以及随机波动率远期利率的鞅条件的精确解；
并给出非线性理论的计价物变换、债券期权价格与远期利率定价核（pricing
kernel）的哈密顿量推导（见第10章）。至此，第三部分在第7章"路径积分侧"与
第10章"哈密顿量侧"各形成一条完整的理论闭环。

\section{体例与使用建议}

原书第1章结尾交代全书体例。各章正文聚焦量子形式应用于金融的概念与理论
方面；更数学性的材料放进各章之后的附录（Appendix）。少数几处读者会受益于
更多细节，推导直接写进正文但用小号字体排印；书末另有一个总附录，收录对
全书材料起辅助作用的数学结果。设置这些附录的用意是把所有专题处理得完整
而全面：打算动手复算的读者会用上它们；原则上，各章附录与小号字推导都可以
跳过，不损失任何连贯性。

本导读增补一条使用建议：初读者把本章当目录，按第2--10章的预告分配深浅；
要动手计算的读者盯住上文那份"范例清单"——第4章障碍期权、第5章随机波动率、
第6章Vasicek与HJM、第7--8章刚性高斯模型——它们是全书精确可解结果的完整
名单。术语与记号可随时查卷末的术语中英对照表与记号总表。

%TAGS: （原书第1章Synopsis无编号公式、无图表，tag清单为空；已对照PNG逐页核验，p019--p022为正文，p023为Part I扉页，p024空白）
